%=====================================================================
\chapter{Generative AI in the Research Process}\label{ch:genai}
%=====================================================================
\locator{5}{Part~5 --- Rigour, Reproducibility and Ethics}

\begin{outcomes}
By the end of this chapter you will be able to:
\begin{tight}
  \item \textbf{Quote} the two clauses of the IUB regulations that govern AI
    use, and explain the tension between them.
  \item \textbf{Classify} a proposed use of a generative model as permitted,
    prohibited or requiring disclosure.
  \item \textbf{Explain} why AI detection tools cannot settle an accusation.
  \item \textbf{Write} an AI-use disclosure statement.
  \item \textbf{Reconcile} the institutional position with international
    publisher policy when you submit abroad.
\end{tight}
\end{outcomes}

\begin{prereq}
\chref{integrity}, in particular the definition of plagiarism and the
similarity thresholds --- machine-generated text is placed inside that
definition by the regulations, which is unusual and consequential.
\end{prereq}

\begin{keyterms}
generative AI $\bullet$ large language model $\bullet$ machine-generated text
$\bullet$ accountability $\bullet$ disclosure $\bullet$ authorship $\bullet$
hallucination $\bullet$ AI detection $\bullet$ false accusation
\end{keyterms}

\section{Two Clauses, Read Together}

The IUB position is stricter than most international publisher policy, and it
is stated in two places that pull in different directions. You need both.

\begin{regbox}[Machine-generated text as plagiarism --- PhD \S32(I)(C)(vi), MPhil/MS \S49(C)(vi)]
Listed among activities that ``still count as academic cheating and plagiarism
if done without permission from the original artists/creators'':

\medskip
``(vi) Using ChatGPT and similar machine-generated text''.
\end{regbox}

\begin{regbox}[Permitted use --- PhD \S32(IV), MPhil/MS \S49(4)]
``Researchers may use ChatGPT, Artificial Intelligence (AI) and AI-assisted
technologies \textbf{to understand basic phenomena} of anything and
\textbf{should not replace the key researcher tasks such as producing
scientific insights, analyzing and interpreting data or drawing scientific
conclusions}. The authors are responsible and accountable for the contents of
the work and \textbf{should not rely solely on AI-generated content}.''
\end{regbox}

\begin{alertbox}[The tension, stated plainly]
The first clause treats using machine-generated text as plagiarism. The second
permits AI use for understanding while prohibiting it for the intellectual core
of research. Taken literally and together, the workable reading is:

\medskip
\textbf{AI may help you understand. It may not produce the intellectual content
of your research, and its output may not appear as text in your submission. You
are accountable for everything either way.}

\medskip
This chapter proceeds on that reading. Where your own department or supervisor
interprets it differently, \emph{their} interpretation governs --- ask, in
writing, and keep the answer. What you must not do is assume a permissive
reading because the clauses are in tension.
\end{alertbox}

\section{A Classification for This Course}

The following is the course's operating rule. It is consistent with the reading
above and is deliberately conservative where the regulations are ambiguous.

\begin{center}
\small
\begin{longtable}{@{}L{5.4cm}L{2.4cm}L{6.0cm}@{}}
\toprule
\textbf{Use} & \textbf{Status} & \textbf{Reasoning} \\
\midrule
\endfirsthead
\toprule
\textbf{Use} & \textbf{Status} & \textbf{Reasoning} \\
\midrule
\endhead
\bottomrule
\endfoot
Explaining a concept you do not understand, as you would use a textbook
  & Permitted
  & Squarely ``to understand basic phenomena''. Verify against a citable
    source \\
Explaining an error message or unfamiliar API
  & Permitted
  & Same category; no research content is produced \\
Suggesting search terms you might have missed
  & Permitted, verify
  & The search itself must still be systematic and reported
    (\chref{slr}) \\
Translating your own sentence into English
  & Disclose
  & The idea is yours; the wording is generated. Disclose and take
    responsibility for accuracy \\
Improving grammar of text \emph{you wrote}
  & Disclose
  & Widely accepted internationally; disclose given \S32(I)(C)(vi) \\
Transcribing your own interview recordings
  & Disclose
  & Mechanical, not interpretive --- but check every transcript against the
    audio, and note the tool (\chref{qualitative}) \\
Boilerplate code you would otherwise copy from documentation
  & Disclose
  & Disclose in the artefact, and you remain responsible for correctness
    (\chref{reproducibility}) \\
\midrule
Drafting any section of the thesis or paper
  & \textbf{Prohibited}
  & Machine-generated text under \S32(I)(C)(vi) \\
Writing the literature review or summarising papers you have not read
  & \textbf{Prohibited}
  & Also breaches \chref{appraisal}: never cite what you have not opened \\
Generating research questions or hypotheses
  & \textbf{Prohibited}
  & ``Producing scientific insights'' \\
Interpreting your results or writing the discussion
  & \textbf{Prohibited}
  & ``Analyzing and interpreting data or drawing scientific conclusions'' \\
Coding qualitative data or generating themes
  & \textbf{Prohibited}
  & Interpretation, explicitly (\chref{qda}) \\
Producing citations or reference lists
  & \textbf{Prohibited}
  & Models fabricate plausible references. See below \\
Generating data, or filling gaps in data
  & \textbf{Prohibited}
  & Fabrication (\chref{integrity}) \\
Reviewing someone else's manuscript
  & \textbf{Prohibited}
  & Confidentiality breach; publishers prohibit it explicitly \\
\end{longtable}
\end{center}

\begin{conceptbox}[The test to apply to a case not in the table]
Ask: \emph{does the output of this use constitute part of the intellectual
contribution, or part of the submitted text?}

\medskip
If either, it is prohibited. If neither --- if it only changed what \emph{you}
understand, and you then produced the contribution and the text yourself --- it
is permitted. When genuinely unsure, ask your supervisor in writing before
doing it, not afterwards.
\end{conceptbox}

\section{Fabricated Citations}

\begin{alertbox}[The failure mode that ends degrees]
Language models generate references that are formatted perfectly, attributed to
real and plausible authors in the right field, published in real journals ---
and that do not exist. Others exist but say something different from what is
claimed.

\medskip
Under \reg{PhD Regs 2024}{\S32(I)(A)(viii)}, ``giving incorrect information
about the source of a quotation'' is plagiarism. A fabricated citation in a
thesis is therefore not an embarrassing slip; it falls inside the plagiarism
definition, and an external examiner who cannot find a cited paper will ask why.

\medskip
The defence is the one \chref{appraisal} already requires: every citation in
your work corresponds to a paper you have opened and read, entered into your
reference manager from the publisher's own metadata (\chref{bibliometrics}).
Follow that and the failure mode cannot occur.
\end{alertbox}

\section{Detection Does Not Work, and That Cuts Both Ways}

\begin{center}
\small
\begin{tabular}{@{}L{3.6cm}L{10.7cm}@{}}
\toprule
False negatives
  & Lightly edited generated text evades detectors reliably. A detector that
    can be defeated by paraphrasing does not enforce anything \\
False positives
  & The serious problem. Detectors flag human writing, and they do so
    unevenly: non-native English writers, and anyone writing in the plain,
    formulaic register that scientific prose rewards, are flagged
    disproportionately \\
No ground truth
  & Unlike the similarity report, which points at a specific matched source you
    can inspect, a detector produces a probability with nothing to check \\
\bottomrule
\end{tabular}
\end{center}

\begin{conceptbox}[What follows for you, and for anyone judging you]
\textbf{If you are accused.} A detector score is not evidence in the way a
similarity match is. What answers an accusation is process evidence: dated
drafts in version control, reading notes in your reference manager, memos,
and the ability to explain your own argument in conversation. Keep the audit
trail (\chref{qda}, \chref{reproducibility}) --- it protects you here as well.

\medskip
\textbf{If you are judging.} Never rest a misconduct finding on a detector
score alone. \reg{PhD Regs 2024}{\S29(V)(ii)} already requires that faculty
``verify each similarity index to identify potential clues'' rather than acting
on a number; the same discipline applies with more force where there is no
matched source to verify against.

\medskip
\textbf{For everyone.} This is precisely why disclosure matters. A stated,
honest account of what you used is worth more than any detection technology,
and it is the only mechanism that actually scales.
\end{conceptbox}

\section{Where the Institutional and International Positions Differ}

You will submit to journals whose policies were written elsewhere. They
converge on three points and diverge from IUB on a fourth.

\begin{center}
\small
\begin{tabular}{@{}L{4.4cm}L{4.6cm}L{4.7cm}@{}}
\toprule
 & \textbf{Common publisher policy} & \textbf{IUB regulations} \\
\midrule
AI as author
  & Prohibited --- AI cannot take responsibility, which is an authorship
    criterion~\cite{icmje2024recommendations}
  & Same; authors ``responsible and accountable'' \\
Accountability
  & Authors fully responsible for all content
  & Same, explicitly \\
AI in peer review
  & Prohibited --- confidentiality
  & Not addressed; treat as prohibited \\
AI-assisted writing
  & Generally \emph{permitted with disclosure} --- language polishing is
    commonly accepted
  & \textbf{Stricter}: machine-generated text is listed as plagiarism \\
\bottomrule
\end{tabular}
\end{center}

\begin{alertbox}[Comply with the stricter rule]
A practice permitted by Elsevier or Springer with a disclosure statement may
still breach \S32(I)(C)(vi) for your thesis. Your degree is governed by the
IUB regulations; the journal's policy governs the journal.

\medskip
Where they differ, follow the stricter one and disclose in both places. That is
always safe. The reverse --- relying on a permissive publisher policy in a
document examined under IUB regulations --- is not.
\end{alertbox}

\section{Disclosure}

\begin{templatebox}[C15 --- AI-use disclosure statement]
Place in the thesis front matter and in the methods section of any paper. If
you used nothing, say that --- an explicit statement is better than silence.

\medskip
\textbf{Tools used.} Name, version or access date.
\filllines{1}

\noindent\textbf{Purpose.} Exactly what each was used for, in specific terms.
``To understand the difference between fixed and random effects before reading
[ref]'' --- not ``for assistance''.
\filllines{2}

\noindent\textbf{Where it was \emph{not} used.} State that no part of the
argument, analysis, interpretation or submitted text was generated.
\filllines{2}

\noindent\textbf{Verification.} How you checked anything the tool told you ---
against which sources.
\filllines{2}

\noindent\textbf{Accountability.} ``I am responsible for all content of this
work, including any part informed by the tools named above.''
\filllines{1}

\medskip
\noindent Signature: \fillline{5cm} \quad Date: \fillline{3cm}
\end{templatebox}

\section{What to Do Instead}

The uses most tempting to a scholar under time pressure are exactly the
prohibited ones. Each has a legitimate substitute that is also better research
practice.

\begin{center}
\small
\begin{tabular}{@{}L{4.4cm}L{9.9cm}@{}}
\toprule
\textbf{The temptation} & \textbf{The substitute} \\
\midrule
``Summarise these 40 papers''
  & The three-pass method (\chref{appraisal}) and a structured extraction form
    (\chref{slr}). Slower, and it produces the synthesis a summary cannot \\
``Write my literature review''
  & An annotated bibliography built as you read, then a synthesis organised by
    theme rather than by paper \\
``Interpret these results''
  & The alignment table (\chref{questions}) --- your interpretation is what
    connects the result back to the question you asked \\
``Improve my English''
  & Your supervisor, a writing centre, a colleague. \chref{writing} is largely
    about writing that does not need rescuing \\
``Generate my research questions''
  & \chref{contribution} and \chref{slr}. A question generated without reading
    the literature cannot have a defensible gap behind it \\
\bottomrule
\end{tabular}
\end{center}

\begin{conceptbox}[The examination is the real constraint]
Set aside the regulations for a moment. You will defend this work in a viva
before examiners who will ask why you chose this design over that one, what
would have falsified your claim, and why the third result came out as it did
(\chref{defending}).

\medskip
Nothing you did not think through yourself is available to you in that room.
A thesis assembled from generated text is not primarily a compliance problem;
it is a defence you cannot give. That is the argument that should settle the
question, and it does not depend on any clause.
\end{conceptbox}

\begin{traditions}{generative AI in research}
\tradEMP{Analysis and interpretation are prohibited uses. The analysis script
must be yours and reviewable, and generated code in an analysis pipeline is a
reproducibility and correctness liability (\chref{reproducibility}).}
\tradDSR{An artefact may legitimately \emph{contain} AI components --- that is
the research. Using AI to design or evaluate the artefact on your behalf is a
different matter and is prohibited.}
\tradQUAL{Coding and theme generation are interpretation, explicitly prohibited
by \S32(IV). The researcher is the instrument (\chref{qualitative}); delegating
it removes the method.}
\tradFORM{Proof assistants and SMT solvers are not generative AI and are long
established. A machine-checked proof is stronger evidence than a human-checked
one --- disclose the tool, as you would any dependency.}
\end{traditions}

\begin{lab}[Write your own AI policy]
\begin{tightnum}
  \item Complete template C15 honestly for work you have already done.
  \item List every use you have made or considered, and classify each against
    the table in this chapter.
  \item For anything you cannot classify confidently, write the question you
    would put to your supervisor.
  \item Send those questions to your supervisor and keep the written answer.
  \item Draft the paragraph you would use to respond to an unfounded accusation
    that your thesis was generated. What evidence do you actually hold? If the
    answer is none, start keeping it now.
\end{tightnum}

\medskip
Step 5 is the one worth doing this week rather than later.
\end{lab}

\begin{checkpoint}
\begin{tightnum}
  \item Quote the two governing clauses and state the reading this chapter
    adopts.
  \item Why is a fabricated citation not merely an error but a breach of
    \S32(I)(A)?
  \item Why can an AI detector score not support a misconduct finding, and what
    can?
  \item A publisher permits AI language polishing with disclosure. May you use
    it in your IUB thesis? Explain.
\end{tightnum}
\end{checkpoint}

\begin{chaptersummary}
\begin{tight}
  \item Two clauses govern: machine-generated text is listed as plagiarism, and
    AI may be used to understand but not to produce insights, analyse or
    interpret. You are accountable regardless.
  \item Working reading: AI may help you understand; it may not produce your
    intellectual contribution or the text you submit.
  \item Permitted with disclosure: translation of your own sentences, grammar
    on your own text, transcription, boilerplate code. Prohibited: drafting,
    literature review, question generation, interpretation, qualitative coding,
    citations, data, peer review.
  \item Fabricated citations fall inside the plagiarism definition. Cite only
    what you have opened.
  \item Detection is unreliable in both directions. Your defence is a dated
    audit trail; disclosure is the mechanism that actually works.
  \item Where publisher policy is more permissive than IUB's, follow IUB's.
  \item The viva is the real constraint: you cannot defend reasoning you did
    not do.
\end{tight}
\end{chaptersummary}

\begin{reviewq}
  \item \S32(I)(C)(vi) and \S32(IV) are in tension. Draft a replacement clause
    that preserves the regulations' evident intent while removing the
    contradiction, and say what your version permits that the current text
    arguably does not.
  \item AI detectors flag non-native English writers disproportionately. What
    obligations does this place on an institution that uses them, and what
    procedural safeguards would you require before any finding?
  \item International policy permits AI language polishing with disclosure;
    IUB's does not. Given that scholars writing in a second language face a
    real disadvantage in publication, is the stricter rule fair? Argue both
    sides and propose a policy.
  \item If a thesis cannot be distinguished from one written with substantial
    AI assistance by reading it, what does that imply about how theses should
    be assessed? Design an assessment that would make the distinction visible.
\end{reviewq}
